%% ma: L12-B14-0047
%% lop: 12
%% chuong: 5
%% bai: 14
%% dang: 12.14.1
%% loai: TN4
%% muc: NB
%% dap_an: "B"
%% nguon: "To_n_12_PT_M_t_Ph_ng.pdf (Đề kiểm tra đánh giá kiến thức, Chuyên đề Phương trình mặt phẳng, Nguyễn Anh Khôi - Lê Minh Khôi), Câu 1"
%% kien_thuc: "Bài 14 (phương trình tổng quát của mặt phẳng, vectơ pháp tuyến)"
%% trang_thai: chinh-thuc
%% ghi_chu: "Tài liệu gốc không kèm đáp án; đáp án do Claude giải và kiểm bằng sympy."
\begin{ex}
Trong không gian $Oxyz$, cho vectơ $\vec n=(0;1;1)$. Mặt phẳng nào trong các mặt phẳng được cho bởi các phương trình dưới đây nhận vectơ $\vec n$ làm vectơ pháp tuyến?
\choice{$x=0$}{\True $y+z=0$}{$z=0$}{$x+y=0$}
\loigiai{Mặt phẳng $Ax+By+Cz+D=0$ có vectơ pháp tuyến $(A;B;C)$. Mặt phẳng $y+z=0$ có vectơ pháp tuyến $(0;1;1)=\vec n$. Chọn B.}
\end{ex}
